\subsection{Weighted Flag Certificates}\label{app:multi-certificates}

We use the coordinates and weighted moment matrices defined in
\textbf{Subsection~\ref{sec:weighted-counts}}.

\subsubsection{Bounds in the Simple Support}\label{app:multi-projection}

Complete the simple support to a $D$-regular simple graph using
blue edges. This preserves the induced cycles with all five edges red
and the parts $U_i$ of each such cycle. With normalization $2sD^{n-2}$,
let $q_C$ and $L_R$ be as in \textbf{Appendix~\ref{app:cycle-estimates}}. Thus $L_R$
weights each ordered cycle by $\frac{R}{D}$, where $R=5D-2|B|$. If
$R\ge\alpha D$ for every such cycle, then $L_R\ge\alpha q_C$. The four
simple certificate inequalities used here are
\begin{equation}\label{eq:support-projection-four}
\begin{aligned}
 S-0.99999q_C
   &\le \frac{5}{4}-\frac{0.99999}{8}+2\cdot10^{-9},\\
 S+0.5682L_R
   &\le \frac{5}{4}+10^{-9},\\
 S+0.4897(L_R-0.32q_C)
   &\le \frac{5}{4}-\frac{0.4897\cdot0.32}{8}+2\cdot10^{-9},\\
 S+0.45119\left(L_R-\frac{q_C}{2}\right)
   &\le \frac{5}{4}-\frac{0.45119}{16}+2\cdot10^{-9}.
\end{aligned}
\end{equation}
Each displayed coefficient is rounded down from the corresponding
rational simple certificate in the reproduction guide. If the
quantity multiplied by that coefficient is nonnegative after
moving the constant term to the left, rounding down preserves the
inequality. If that quantity is negative, the rounded inequality
follows from $S\le\frac{5}{4}$.

Substitute $L_R\ge\alpha q_C$ into each inequality. Then $S$ is
bounded by the minimum of five affine functions of $q_C$. The
maximum of this minimum over $q_C\ge0$ is attained at zero or at a
nonnegative intersection of two of the lines, since the line from
the second displayed inequality has negative slope for $\alpha>0$.
The program evaluates these intersections as rationals. It may
also use the independent support bounds $S\le\frac{32}{27}$,
$\frac{145}{121}$ or $\frac{204}{169}$ for $D=9,11,13$, respectively.
To express these inequalities in the weighted model, multiply them
by $1-e=\frac{s}{m}$. This changes the normalization of the cycle
count and its $R$-weighted sum from $2sD^{n-2}$ to $2mD^{n-2}$.
The cycles and their values of $R$ are the same in either completion.
The resulting bound $S\le S_{\max}$ gives $e\ge1-\frac{S_{\max}}{M}$.

\subsubsection{Degree Constraints at Complete Red Subgraphs}

Let $K$ be the vertex set of a red $K_k$ in the weighted host.
For each vertex $v$ put $a_v=|N_{G_{\rm s}}(v)\cap K|$ and let
$h_v$ and $d_{\rm blue}(v)$ count the marked and the unmarked edges
incident with $v$, respectively. Let $m_K$ count the marked edges with
both endpoints in $K$, with multiplicity, and put
$J_K=\sum_{v\in K}j_v$. For any marked
edge, the union of its endpoints' neighborhoods in $K$ meets
every edge of the red clique. In particular, its endpoint
$a$-values sum to at least $k-1$. They sum to at least $k$ if
one endpoint belongs to $K$, and to $2(k-1)$ if both do. Hence
\[
 (k-1)m+\sum_{v\in K}h_v+(k-3)m_K
 \le\sum_v a_vh_v
 \le D(kD-J_K)-\sum_{v\in K}d_{\rm blue}(v).
\]
For the last inequality, use
$h_v=D-d_{\rm blue}(v)$ and
$a_vd_{\rm blue}(v)\ge |N_{\rm blue}(v)\cap K|$.
It follows that
\[
 (k-1)m\le kD(D-1)-DJ_K-(k-3)m_K.
\]
For $k=5$, the right side is at most $5D(D-1)-20$. This
excludes a red $K_5$ both for $M\ge\frac{5}{4}$ and for the odd
counterexample in \eqref{eq:multi-mass}. When $D\le23$ in
the odd reduction, Appendix~\ref{app:multi-no-k4} excludes
a red $K_4$ as well. The resulting forbidden patterns have
zero coordinates. Red triangles may occur.

These exclusions also give inequalities on common red neighborhoods.
For example, if there is no red $K_4$, the red neighborhood of a
vertex contains no red triangle. For any nonnegative weights $w_x$
on that neighborhood,
\[
 2\sum_{xy\text{ ordered red edge}}w_xw_y
 \le \left(\sum_x w_x\right)^2.
\]
To see this, move the combined weight of two nonadjacent vertices
onto the one with greater weighted neighbor sum. The edge sum does
not decrease, and iterating reduces the support of the weights to
a clique of size at most two. The inequality then follows from
$4w_xw_y\le(w_x+w_y)^2$. The weighted calculation uses
$w_x=1,b_{vx},1+b_{vx}$ for a red neighborhood rooted at $v$.
The square on the right includes two choices of the same vertex.

\subsubsection{Additional Linear Inequalities and Exact Duals}

Fix a red pair $xy$. Let $U=N_{G_{\rm s}}(x)\setminus N_{G_{\rm s}}(y)$
and let $W$ be the vertices adjacent to neither $x$ nor $y$. In the
host induced on $U\cup W$, mark only the previously marked edges
between $U$ and $W$. Edges within a part can still join endpoints of
these marked edges. For each marked edge $uv$ between $U$ and $W$, and
each red neighbor $z$ of $y$, the strong clique condition for $uv$
and $yz$ forces $uz$ or $vz$ to be a host edge. If $uv$ meets $y$,
the edge $yz$ itself gives the adjacency. All such vertices $z$ lie
outside $U\cup W$. Hence the Ore-degree $\sigma$ of the bipartite
marked subgraph in this induced host is at most $2D$ minus the
number of red neighbors of $y$.

We use three bounds with this Ore-degree. The total number of marked
edges is at most $\frac{2973}{10000}\sigma^2$, by the companion proof
on witness deficits and neighborhood differences. Its last step is the
rational interval check in \textbf{Appendix~\ref{app:ore-replay}}. The base copies
form a strong clique in the simple support, where the Ore-degree can
only decrease, so the simple rational Ore bound gives at most
$\frac{273323}{10^6}\sigma^2$ base copies. The bound
$\frac{1}{4}\sigma^2$ on the extra copies comes from reserving the base
copies as witnesses, as proved in the companion note on that bound.
The proofs of the triangle and five-cycle inequalities and of the
other constraints for bounded degree are indexed in
\path{reproducibility/multigraph/REPLAY.md}.

A pointwise inequality at a fixed root gives a linear inequality for
the weighted counts. Multiply by the product of the extra multiplicities
on the specified forest edges of the type and sum over all maps of
the type. Local degree terms and counts of
individual neighbors use one more position. Quadratic terms and
counts of an additional edge use two. At $D=9$, for example, the
external neighborhood union at a double pair has size $13$ or
$14$. In a flag extension for the entire neighborhood union, the
images of the two root positions are also valid choices. Thus the
corresponding interval is $[15,16]$. At $D=13$ and $15$ we also
enumerate the possible integer counts of outside vertices adjacent
to two or three vertices of a red triangle, together with the
degree and multiplicity losses at its vertices. In the host on the
marked endpoints before regularization, let $E$ be the sum of the
degree deficits and extra marked degrees at the triangle. Let $t$
count outside vertices with at least two distinct triangle neighbors,
and $q$ those with three. The number of vertices with a triangle
neighbor, including the roots, is $N_+=3(D-1)-E-t-q$. A marked edge
meeting a vertex without a triangle neighbor has its other endpoint
among the $t$ vertices. Thus the number $m_{\rm out}$ of those edges
is at most $(D-2)t-q$. The marked degree sum at the $N_+$ vertices
is $2m-m_{\rm out}$, so
\[
 DE+2t+(D+1)q\le 3D(D-1)-2m.
\]
The generator uses this inequality, degree bounds at the outside
vertices with a unique triangle neighbor, and the finite root-pair
bounds. Filling unused degree with blue edges can only decrease $E$
and preserves $t,q$. The supporting planes of the downward closure
in $E$ of the finite outer bound give linear inequalities in the
rooted counts. They are generated in integer arithmetic by
\path{finite_triangle_profiles.py}. When summing extensions for
$t$, the three maps of the fourth position to a root are subtracted.

Write $AX=(1,0,\ldots,0)^{\mathsf T}$ for the degree and sampling
equations, $NX\le0$ for the linear inequalities, and $G_i(X)\succeq0$
for the matrix inequalities. A stored bound consists of a vector
$\lambda$, a nonnegative vector $\rho$, and matrices $P_i$
given as nonnegative sums of rational rank-one outer products.
Let $G_i^*$ be the adjoint defined by
$(G_i^*(P_i))^{\mathsf T}X=\tr(P_iG_i(X))$. For each feasible $X$,
\[
 \left(A^{\mathsf T}\lambda+N^{\mathsf T}\rho
        -\sum_iG_i^*(P_i)\right)^{\mathsf T}X
 \le\lambda_0 .
\]
Let $r$ be the coefficient vector on the left minus the coefficient
vector of $M$. If $0\le X_j\le U_j$ for each coordinate $j$, then
$M\le\lambda_0+\sum_{j:r_j<0}(-r_j)U_j$. For a coordinate $X(P,T)$,
an upper bound is
$\frac{1-e}{|\operatorname{Aut}(P,T)|}$ for $T=\varnothing$ and
$\frac{e}{|\operatorname{Aut}(P,T)|}$ otherwise. To obtain it, start at
a red edge, in $T$ if $T$ is nonempty, and extend along a
spanning tree containing $T$. Each new position contributes at
most $D$ in support count or extra multiplicity. The replayers
also use bounds on sums of coordinates of equal order and
equal forest size, obtained by counting all labeled spanning
trees containing the specified red root. Positive residuals
can be dropped since the coordinates are nonnegative. The
total compensation is bounded at the endpoints of the allowed
interval for $e$. For an infeasibility certificate, the objective
is zero and the compensated right side is negative.

\subsubsection{Verification of the Quadratic and Odd-Degree Bounds}

The five-position certificates first give $M<1.28746$ and
$e<\frac{421}{2000}$ when $M>\frac{5}{4}$. The
certificate at $M\ge\frac{4967}{4000}$ gives the second claim of
\textbf{Lemma~\ref{prop:multi-uniform}}. Under the assumption $M>\frac{5}{4}$,
the strengthened five-position certificate
\path{direct_half_scs.json} improves the mass bound to $1.28492$.
In the six-position chain, the successive lower ends of the
assumed mass intervals are $1.283$, $1.279$, $1.274$, $1.268$
and $\frac{5}{4}$. Put $r=\frac{R}{D}$, $j=\frac{\sum_i j_{c_i}}{D}$ and
$g=M-\frac{5}{4}$. The count in \textbf{Appendix~\ref{app:multi-structure}} gives
$r>\frac{1}{2}$ when $g>0$. For $\frac{1}{2}<r\le1$, substituting
$K\le R-\sum_i j_{c_i}$ in \eqref{eq:multi-rough-discharge}
gives $2g\le r(r-\frac{1}{2})-rj$.
Dropping the nonnegative loss term in \eqref{eq:multi-joint}
gives the stronger bound
\[
 2g\le \frac{1+r}{2+r}r\left(r-\frac{1}{2}\right)-\frac{1+r^2}{2+r}j .
\]
The replay uses the relevant inequality with $j\ge0$ and the
current lower bound on $M$ to obtain a lower bound on $r$. All
chosen lower bounds are at most one, so a cycle with $r>1$
satisfies them as well. The simple support bound then supplies
the lower end of the interval for $e$. The replay verifies the
exact comparisons between consecutive intervals. In the last
interval the dual gives
\[
 M\le\frac{124968062205915704703985647}
               {100000000000000000000000000}<\frac{5}{4}.
\]
This contradicts the assumed strict violation and proves
\textbf{Lemma~\ref{prop:multi-uniform}}.

\begin{lemma}\label{lem:multi-even-equality}
Suppose the host has positive even maximum degree $D$ and the
number of marked copies is $\frac{5D^2}{4}$. The host induced by
the marked endpoints is a blow-up of $C_5$ with all five parts of
size $\frac{D}{2}$. Every edge of that blow-up is single and marked.
\end{lemma}

\begin{proof}
The five-position multiplicity bound applies at $M=\frac{5}{4}$ and
gives $e<0.2105$. If every induced red five-cycle had
$\frac{R}{D}\ge\frac{1}{2}$, the support bound in
\textbf{Appendix~\ref{app:multi-projection}} would give $S\le1.22181$ and
hence $e>0.0225$. These bounds satisfy the hypotheses of the last
six-position certificate, whose upper bound is strictly less than
$\frac{5}{4}$. Thus there is a cycle with $R<\frac{D}{2}$.

Apply the redistribution leading to \eqref{eq:multi-rough-discharge}
at this cycle. For every outside vertex $x$ with $k_x>0$, its
resulting contribution is strictly less than $\frac{k_xD}{2}$.
For $k_x=1$, a vertex either keeps at most $R-2$ or sends half of
its marked degree and has a blue anchor edge. For $k_x=2$, there
is a blue anchor edge and no incoming transfer. For $k_x=3$ with
a marked neighbor of signature size zero, the incoming transfer
is at most $\max\{\frac{D+R}{4},R\}<\frac{D}{2}$. Without such a
neighbor, it is at most $\frac{D-3}{2}$, since the three anchor
adjacencies use degree units inside the core. For $k_x\ge4$,
the transfer is at most $D-k_x$ and its own contribution is
at most $D$. Vertices of signature size zero retain no positive
contribution. Since $K\le R$, the presence of an outside vertex
of positive signature size would give $2m<D|B|+\frac{DK}{2}
\le\frac{5D^2}{2}$ by \eqref{eq:multi-core-cut}.

Hence $K=0$. Now $2m\le D|B|=\frac{D(5D-R)}{2}$, so $R=0$.
By \eqref{eq:multi-core-identities}, the five part sizes are
$\frac{D}{2}$. For a vertex of $U_i$, the neighbors in
$U_{i-1}\cup U_{i+1}$ already use all $D$ degree units. There are
$\frac{5D^2}{4}$ edges on these pairs. Any marked copy outside
$B$ would require an outside vertex of positive signature size.
Thus all edges in the blow-up are marked, and $B$ contains every
marked endpoint.
\end{proof}

We next verify the bound used to force an induced red five-cycle in
\textbf{Lemma~\ref{lem:counterexample-cycle}}. For $m,D>0$ we claim that
\begin{equation}\label{eq:absence-red-cycle}
 q_C=0 \quad\Longrightarrow\quad M<\frac{98}{81}.
\end{equation}
Suppose that $q_C=0$ and $M\geq\frac{98}{81}$. The uniform bounds for
the multigraph and its simple support give $M,S\leq\frac54$.
If $b=0$ then $e=0$. If $b>0$, \textbf{Lemma~\ref{lem:multi-extra}}
gives $e\leq1-\frac{M}{2}$. Thus in either case
\[
 \frac{98}{81}\leq M\leq\frac{12883}{10000},
 \qquad 0\leq e\leq\frac{32}{81}.
\]
These are the intervals used in the stored calculation.

Let $P$ be the five-position pattern with five red cycle edges and
color $0$ on all other pairs. A color-preserving map of $P$ is
injective. Adjacent positions must have distinct images, and any
two nonadjacent positions are distinguished by a third position
adjacent to exactly one of them. Hence $q_C=0$ excludes every such map. It follows from
the definition of $X(P,T)$ in
\textbf{Subsection~\ref{sec:weighted-counts}} that $X(P,T)=0$ for
each forest $T\subseteq E(P)$ used in the calculation, including
$T=\varnothing$. Since the coordinates are nonnegative, the verifier
imposes all these equations by summing the distinct coordinates
$X(P,T)$ and requiring this sum to be at most zero.

Together with this inequality, the five-position calculation uses
the weighted degree and flag constraints, the simple-support bound,
the first inequality in \eqref{eq:support-projection-four}, and the
rooted five-cycle bounds indexed in \path{PROOF_NOTES.md}.
The rational coefficients in \path{absence_c5_mass_98_81.json}
give the bound below after the correction for negative residual
coefficients described above.
\[
 M\leq\frac{14802480264533669793521}{12250000000000000000000}
   <\frac{98}{81}.
\]
This contradicts the assumed range of $M$ and proves
\eqref{eq:absence-red-cycle}. The replay command is given in
\textbf{Appendix~\ref{app:multi-reproduction}}.

For the eleven rows of \textbf{Table~\ref{tab:multi-odd}}, the replay
rebuilds the linear systems and checks their projection inputs,
multipliers and residual bounds in rational arithmetic.
\textbf{Appendix~\ref{app:multi-reproduction}} lists the inputs and commands.
